% \iffalse meta-comment
%
% Copyright (C) 2017- by Yanshuo Chu <yanshuoc@gmail.com>
% Distributed under the LaTeX Project Public License, version 1.3a or later.
%
% \fi
%
% ^^A 用户指南；只说明当前可用接口。实现说明由 driver 的开发者开关控制。
% \hitsetup{
%   openright=false,
%   info={title={\hithesis 用户手册}[\hithesis\ User Guide],author={\hithesis 维护组}},
%   overrides={
%     base-document-type={用户手册},
%     base-degree-level={},
%     header={document-type={}},
%     header={university={\hithesis 用户手册},degree-level={}},
%   },
% }
% \cover[degree-type-line=false]
% \tableofcontents[lang={zh,en},depth=2]
%
% \chapter{开始使用}[Getting Started]
% \label{sec:installation}
% \hithesis\ 是哈尔滨工业大学的 \hologo{LaTeX} 模板，支持三个校区的本科、硕士、博士学位论文，开题、中期报告和博士后报告。
% 先从 \file{examples/demo/} 复制示例，再修改类选项、个人信息和正文。
% \hithesis\ 依据学校公开指南与范例实现，提交前请核对院系当年的要求。
%
% \section{安装与编译}[Installation and Compilation]
% 推荐使用较新的完整 \hologo{TeX}~Live，优先选择 \hologo{LuaLaTeX}；\hologo{XeLaTeX}、\hologo{pdfLaTeX} 也支持，
% 但字体、断行和行高可能不同。需要贴近 Word 范例时使用 \hologo{LuaLaTeX} 和学校指定字体。
%
% 模板在 \hologo{LuaLaTeX} 下使用 \pkg{luatexja}，在 \hologo{XeLaTeX} 下使用
% \pkg{xeCJK}；\hologo{pdfLaTeX} 的中文支持由 \cls{ctexbook} 和 \pkg{CJK} 提供。
% \pkg{xeCJK} 所需的 ctex 辅助宏包也会一并加载。封面、章节标题和正文版式由模板与
% \pkg{docxlike} 设置，通常无需另行加载中文宏包。更新发行版或字体后，建议检查文档的
% 换行和分页；论文定稿时，可保留所用的模板、字体及编译环境，方便日后重新编译。
% 编译项目应保留示例附带的 \file{latexmkrc}，它会处理参考文献、成果库和索引。
% \begin{shell}
% latexmk -lualatex final.tex
% latexmk -lualatex report.tex
% \end{shell}
% 切换引擎时先清理辅助文件，避免沿用另一引擎的结果。
% 从源码仓库使用 \hithesis，先在仓库根目录运行：
% \begin{shell}
% make cls
% make distribute
% \end{shell}
% 发布包中已有生成文件时，不必自行解压源码。
%
% \section{最小学位论文}[A Minimal Thesis]
% \label{sec:usage}
% 下例不依赖外部图片或文献库，可直接保存为 \file{main.tex} 编译。
% \begin{latex}
% \documentclass{hithesis}
% \hitsetup{info={
%   title={论文题目}[Thesis Title],author={张三}[Zhang San],
%   supervisor={李四教授}[Prof. Li Si],
%   affiliation={某学院},speciality={某专业},degree-applied={工学博士}[Doctor of Engineering],
%   student-id={1234567890},
% }}
% \begin{document}
% \cover
% \titlepage
% \begin{abstract}[keywords={关键词一,关键词二}]
% 中文摘要。
% \end{abstract}
% \begin{abstract}[lang=en,keywords={keyword one,keyword two}]
% English abstract.
% \end{abstract}
% \tableofcontents
% \chapter{绪论}[Introduction]
% 正文。
% \begin{conclusions}结论。\end{conclusions}
% \begin{acknowledgements}致谢。\end{acknowledgements}
% \end{document}
% \end{latex}
% 实际论文推荐在主文件中用 \cs{input} 或 \cs{include} 组织各章，文件布局参见示例。
%
% \section{选择论文或报告类型}[Choosing the Document Type]
% \label{sec:option}
% 类选项写在 \cs{documentclass} 的方括号里。省略时，\hithesis\ 默认排哈尔滨校区的中文博士学位论文。
%
% \DescribeOption{stage}
% 文档类型，默认 \texttt{final}：学位论文、实践成果或博士后报告。
% \texttt{proposal} 为开题报告，\texttt{interim} 为中期报告。
%
% \DescribeOption{degree-level}
% 学位层次，默认 \texttt{doctor}（博士）；其他取值为 \texttt{bachelor}（本科）、
% \texttt{master}（硕士）、\texttt{postdoc}（博士后研究工作报告）。
% 开题和中期报告省略此项时也按博士处理；博士后固定使用终稿阶段。
%
% \DescribeOption{campus}
% 校区，默认 \texttt{harbin}（哈尔滨）；\texttt{weihai} 为威海，\texttt{shenzhen} 为深圳。
% 本科终稿三校区使用统一版式，报告仍按校区选择。
%
% \DescribeOption{lang}
% 正文语言，默认 \texttt{zh}（中文）；\texttt{en} 为英文。
%
% \DescribeOption{degree-type}
% 学位类型，默认 \texttt{academic}（学术学位）；\texttt{professional} 为专业学位。
%
% \DescribeOption{category}
% 学科门类，默认 \texttt{stem}（理工类）；\texttt{hass} 为人文社科类，\texttt{hss} 是其别名。
%
% \DescribeOption{form}
% 论文形式，默认 \texttt{dissertation}（论文），别名 \texttt{thesis}。
% \texttt{practice} 为实践成果，本科为毕业设计，别名 \texttt{design}。
%
% 报告示例：
% \begin{latex}
% \documentclass[stage=proposal,degree-level=master,campus=harbin]{hithesis}
% \end{latex}
% \hithesis\ 已提供常用数学、图表、列表和定理支持。额外使用 Ti\emph{k}Z 等工具时自行加载；
% \file{hithesis.sty} 是可选的示例辅助宏包，不是文档类本身。
%
% 深圳研究生报告的中英文开题、中期按深圳教务部 2018 版 Word 模板设置。
% 开题封面落款为“深圳校区教务部 制”，中期没有落款；英文版使用英文文种及字段标签。
% 目录与正文分别从阿拉伯数字 1 编号，中期目录不印页眉，正文各节不强制另起一页。
% 开题图表与公式连续编号，中期按节编号；可用编号选项覆盖。
% 博士暂时沿用同阶段、同语种的硕士版式，仅文种随学位变化，不代表已核实博士专用格式。
% 旧博士中期的表格式封面、说明页、正文边框与特殊标题横线不再使用。
%
% \subsection{字体与数学设置}[Fonts and Mathematics]
% 字体预设由模板自己的角色表定义，不再转交给 \pkg{ctex} 的字体预设。中文、西文、数学可分别覆盖。
%
% \DescribeOption{fontset}
% 默认 \texttt{auto}，每个角色独立选择：先尝试 \texttt{windows} 对应字体，
% macOS 上缺项再尝试 \texttt{macos} 对应字体，最后使用 \texttt{fandol} 对应字体。
% 其他平台跳过 macOS 候选。发现某一个 Windows 字体不会切换整套预设。
% 所有显式预设也逐角色检查可用性，缺项回退到 \texttt{fandol}，并在日志中报告。
% 下表以各角色的实际字体排印相同字样，并沿用默认基线校正；编译手册时缺少的字体会注明回退名称。
%
% \par\noindent
% \begin{talltblr}[caption={字体预设的中文角色},label={tab:font-presets-cjk}]
%   {colspec={l X[l] X[l] X[l] X[l]},cells={font=\small},row{1}={font=\small\bfseries}}
% \toprule
% 预设 & 宋体 & 黑体 & 楷体 & 仿宋 \\ \midrule
% windows & \hitFontSample{windows}{songti}{SimSun} & \hitFontSample{windows}{heiti}{SimHei}
%   & \hitFontSample{windows}{kaishu}{KaiTi} & \hitFontSample{windows}{fangsong}{FangSong} \\
% macos & \hitFontSample{macos}{songti}{Songti SC} & \hitFontSample{macos}{heiti}{Heiti SC}
%   & \hitFontSample{macos}{kaishu}{Kaiti SC} & \hitFontSample{macos}{fangsong}{STFangsong} \\
% founder & \hitFontSample{founder}{songti}{FZShuSong-Z01} & \hitFontSample{founder}{heiti}{FZHei-B01}
%   & \hitFontSample{founder}{kaishu}{FZKai-Z03} & \hitFontSample{founder}{fangsong}{FZFangSong-Z02} \\
% noto & \hitFontSample{noto}{songti}{Noto Serif CJK SC} & \hitFontSample{noto}{heiti}{Noto Sans CJK SC}
%   & \hitFontSample{noto}{kaishu}{AR PL KaitiM GB} & \hitFontSample{noto}{fangsong}{朱雀仿宋} \\
% fandol & \hitFontSample{fandol}{songti}{FandolSong} & \hitFontSample{fandol}{heiti}{FandolHei}
%   & \hitFontSample{fandol}{kaishu}{FandolKai} & \hitFontSample{fandol}{fangsong}{FandolFang} \\
% \bottomrule
% \end{talltblr}\par\noindent
% \begin{talltblr}[caption={字体预设的西文角色},label={tab:font-presets-latin}]
%   {colspec={l X[l] X[l] X[l]},cells={font=\small},row{1}={font=\small\bfseries}}
% \toprule
% 预设 & 衬线体 & 无衬线体 & 等宽体 \\ \midrule
% windows & \hitFontSample{windows}{serif}{Times New Roman} & \hitFontSample{windows}{sans}{Arial}
%   & \hitFontSample{windows}{mono}{Consolas} \\
% macos & \hitFontSample{macos}{serif}{Times New Roman} & \hitFontSample{macos}{sans}{Arial}
%   & \hitFontSample{macos}{mono}{Menlo} \\
% founder & \hitFontSample{founder}{serif}{Times New Roman} & \hitFontSample{founder}{sans}{Arial}
%   & \hitFontSample{founder}{mono}{DejaVu Sans Mono} \\
% noto & \hitFontSample{noto}{serif}{TeX Gyre Termes} & \hitFontSample{noto}{sans}{TeX Gyre Heros}
%   & \hitFontSample{noto}{mono}{DejaVu Sans Mono} \\
% fandol & \hitFontSample{fandol}{serif}{TeX Gyre Termes} & \hitFontSample{fandol}{sans}{TeX Gyre Heros}
%   & \hitFontSample{fandol}{mono}{Latin Modern Mono} \\
% \bottomrule
% \end{talltblr}\par
% 数学字体在所有预设、所有平台下均默认使用 TeX Gyre Termes Math；不随中文或西文选择变化。
% macOS 的 Songti 使用 Regular/Bold，Heiti 使用 Light/Medium，Kaiti 使用 Regular/Bold。
% \hologo{LuaLaTeX} 会查找 Apple 字体资源目录中的楷体和仿宋文件。
% \texttt{mac/macnew/macold} 是 \texttt{macos} 的兼容名，\texttt{ubuntu} 对应 \texttt{noto}。
% \texttt{adobe/hanyi} 保留为兼容预设，\texttt{none} 保留为不设置中文字体的兼容值。
% 最终回退只使用 \hologo{TeX}~Live 字体，所需 \texttt{fandol}、\texttt{tex-gyre}、
% \texttt{lm}、\texttt{tex-gyre-math} 已列入项目依赖；最小安装须补齐这些包。
% 朱雀仿宋可由 \texttt{lxgw-fonts} 包提供，文件名为 \texttt{LXGWZhuqueFangsong-Regular.ttf}。
% \texttt{noto} 的楷体采用文鼎 PL 简中楷（AR PL KaitiM GB），
% 由 \texttt{arphic-ttf} 包提供，按文件 \texttt{gkai00mp.ttf} 加载；缺失时回退 FandolKai。
% Noto CJK 与方正字体须另行安装；\hologo{TeX}~Live 的 \texttt{noto} 包不提供 Noto 中文字体。
% DejaVu Sans Mono 由 \texttt{dejavu} 包提供（属于 \texttt{collection-fontsextra}），
% 并非最小安装必备；缺失时回退 Latin Modern Mono。
%
% 也可写角色字典：\texttt{preset} 选择基础预设，未写的角色继承它。
% 角色为 \texttt{songti/heiti/kaishu/fangsong/kaishu-gb2312/lishu/xinwei}、
% \texttt{serif/sans/mono/math}；显式角色覆盖优先于中文、西文分项预设，角色取 \texttt{auto} 时恢复继承。
% 显式字体名不可用时报错，不静默替换。隶书、新魏没有 Fandol 对应字形，封面按 \option{cover-script} 使用轮廓替代。
% \begin{latex}
% \documentclass[fontset={preset=windows,mono=Menlo,fangsong=FangSong}]{hithesis}
% \end{latex}
% 配置表见 \file{src/config/hit-font-presets.dtx}。实际字体与 Word 目标度量分开，换字体不等于改变目标行盒。
% 代码的中文族优先使用仿宋角色；该字体不可用时依次尝试楷体、宋体，并报告实际选用的字体。
% 预设名称相同不能保证 Word 中字体文件也相同，对照时须核对实际字体及版本。
%
% \DescribeOption{fontset/baseline-correction}
% 中文字体视觉基线校正，是 \option{fontset} 的子键，取 \texttt{auto/true/false}，默认 \texttt{auto}。
% \begin{latex}
% \documentclass[fontset={preset=fandol,baseline-correction=false}]{hithesis}
% \end{latex}
% 文档中可用 \cs{hitsetup}\texttt{\{fontset=\{baseline-correction=auto\}\}} 调整校正开关，
% 并支持分组恢复；字体预设和角色仍在类选项中指定。
% \hologo{LuaLaTeX} 下自动启用：测量各角色常规字面中“国”字的墨迹中心，
% 与 Windows 对应字体的参考值比较，将差值按实际字号换算为字形上下偏移。
% 缺少“国”字时尝试“國”；两字都缺失或无法读取轮廓边界时跳过并记录原因。
% 差值小于 0.001\,em 时不移动。粗体、斜体继承常规字面的偏移；中文代码按校正后的实际字号换算。
% 此项不改变字宽、行距和行盒尺寸，也不移动西文及数学字体。
% 自定义字体没有绑定模板角色时不自动校正；\texttt{windows-local/mac-word} 路径方案暂不绑定角色。
% \hologo{XeLaTeX}、\hologo{pdfLaTeX} 暂不支持此校正，\texttt{auto} 不启用，显式 \texttt{true} 给出提示。
% 表中灰线为排版基线，三个字样均为 12\,bp；右列使用 Windows 对应角色作为参照。
% \par\noindent
% \begin{talltblr}[caption={中文字体基线校正对照},label={tab:font-baseline}]
%   {colspec={X[l] c c c},cells={font=\small},row{1}={font=\small\bfseries}}
% \toprule
% 预设／角色 & 未校正 & 已校正 & Windows 参照 \\ \midrule
% fandol／宋体 & \hitBaselineSample{fandol}{songti}{false} & \hitBaselineSample{fandol}{songti}{auto}
%   & \hitBaselineSample{windows}{songti}{false} \\
% fandol／黑体 & \hitBaselineSample{fandol}{heiti}{false} & \hitBaselineSample{fandol}{heiti}{auto}
%   & \hitBaselineSample{windows}{heiti}{false} \\
% fandol／楷体 & \hitBaselineSample{fandol}{kaishu}{false} & \hitBaselineSample{fandol}{kaishu}{auto}
%   & \hitBaselineSample{windows}{kaishu}{false} \\
% fandol／仿宋 & \hitBaselineSample{fandol}{fangsong}{false} & \hitBaselineSample{fandol}{fangsong}{auto}
%   & \hitBaselineSample{windows}{fangsong}{false} \\
% \bottomrule
% \end{talltblr}\par
% Windows 参照字体未安装时，右列显示该角色的回退字体，不能作为 Windows 字形对照。
%
% \DescribeOption{fontset-zh}
% 单独指定中文字体，默认 \texttt{auto}，跟随 \option{fontset} 或自动选择。
% 除上述中文方案外，还可取 \texttt{windows-local}（从指定目录读中易字体）、
% \texttt{mac-word}（读取 macOS 版 Word 附带字体）。
%
% \DescribeOption{fontset-en}
% 西文字体，默认 \texttt{auto}，继承整套预设；可单独选择上述标准预设的西文角色。
% 也接受 \texttt{times}（Times New Roman、Arial、Consolas）、
% \texttt{termes}（\hologo{TeX} Gyre Termes、Heros、Cursor）、\texttt{none}（不设置）。
% 旧 \texttt{times} 方案在 Times 缺失时回退到 \texttt{termes}；预设缺项按角色回退。
% \hologo{LuaLaTeX} 和 \hologo{XeLaTeX} 下，等宽字体以 Consolas 为字宽目标，
% 按正文字号校正并取最近的半磅，其他字号沿用同一比例。12\,bp 正文中的 Consolas 为 12\,bp，
% 替代字体可能使用较小字号；\texttt{fontset-en=none} 时由用户自行设置。
%
% \DescribeOption{fontset-math}
% 数学字体，默认 \texttt{auto}；\hologo{LuaLaTeX}、\hologo{XeLaTeX} 默认使用 TeX Gyre Termes Math，
% 仅显式 \texttt{math} 角色或本选项可以覆盖，
% \hologo{pdfLaTeX} 默认使用 newtxmath。
% 可取 \texttt{cambria}（Cambria Math）、\texttt{termes}（\hologo{TeX} Gyre Termes Math）、\texttt{xits}（XITS Math）、
% \texttt{stix}（STIX Two Math）、\texttt{libertinus}（Libertinus Math）、
% \texttt{newcm}（New Computer Modern Math）、\texttt{latinmodern}（Latin Modern Math）、
% \texttt{newtx}（newtxmath）、\texttt{none}（不设置）。
%
% \DescribeOption{math-style}
% 数学字母正斜体规则，默认 \texttt{auto}，按 GB 处理。
% \texttt{TeX} 令大写希腊字母正体、小写斜体；\texttt{ISO} 令希腊字母斜体；
% \texttt{GB} 在 ISO 基础上将偏导符号改为正体。
% 此项由 \pkg{unicode-math} 实现，\texttt{newtx} 方案使用宏包自身的规则。
%
% \DescribeOption{windows-font-dir}
% 中易字体文件目录，默认当前目录；供 \texttt{windows-local} 和 \hologo{pdfLaTeX} 查找字体使用。
%
% \DescribeOption{cover-script}
% 封面美术字的绘制方式，默认 \texttt{auto}：有字体则用字体，否则绘制轮廓。
% \texttt{font} 优先使用字体，缺失时提示并回退；\texttt{outline} 直接绘制轮廓。
%
% \subsection{别名与旧选项}[Aliases and Legacy Options]
% 以下写法仍可识别，新文档建议使用前文的名称。
%
% \DescribeOption{type}
% \option{degree-level} 的旧名，接受同一组学位取值，使用时给迁移提示。
%
% \DescribeOption{language}
% \option{lang} 的旧名；\texttt{chinese/english} 对应 \texttt{zh/en}，也接受 \texttt{zh/en}，使用时给迁移提示。
%
% \DescribeOption{stage}
% 旧值 \texttt{opening} 对应 \texttt{proposal}，\texttt{midterm} 对应 \texttt{interim}，两者均给迁移提示。
% 历史拼写 \texttt{mid-term} 也对应 \texttt{interim}。
%
% \DescribeOption{category}
% \texttt{hss} 与 \texttt{hass} 同义，均表示人文社科类。
%
% \DescribeOption{form}
% \texttt{thesis} 与 \texttt{dissertation} 同义；\texttt{design} 与 \texttt{practice} 同义。
%
% \DescribeOption{font-en}
% 与 \option{fontset-en} 同义；旧名 \option{font} 也接受，但会给迁移提示。
%
% \DescribeOption{font-zh}
% 与 \option{fontset-zh} 同义；旧名 \option{cjk-font} 也接受，但会给迁移提示。
%
% \DescribeOption{math-font}
% 与 \option{fontset-math} 同义，选择数学字体方案。
%
% \DescribeOption{design}
% 旧布尔选项：\texttt{true} 对应 \texttt{form=practice}，\texttt{false} 对应论文；使用时给迁移提示。
%
% \DescribeOption{newtxmath}
% 旧数学字体开关：\texttt{true} 在数学字体为 auto 时选择 newtxmath，当前建议写 \texttt{fontset-math=newtx}。
%
% \texttt{degree-level=bachelor*} 已不再支持，应写 \texttt{degree-level=bachelor,form=practice}。
% 排版开关的旧类选项见第~\ref{sec:upgrade}~节，当前应通过 \cs{hitsetup} 或部件参数设置。
%
% \chapter{配置与文档结构}[Configuration and Document Structure]
% \section{配置入口}[Configuration Overview]
% \label{sec:keyfamily}
% \DescribeMacro{\hitsetup}
% \cs{hitsetup}\marg{键值表} 设置全文配置，后写的值覆盖先写的值，通常放在导言区。
%
% \DescribeOption{info}
% 设置题目、作者、导师、日期等元信息，字段说明见第~\ref{sec:titlepage}~节。
%
% \DescribeOption{overrides}
% 修改封面标签、部件标题等固定文字，按页面分组设置。
%
% \DescribeOption{styles}
% 设置字体、字号、行距和缩进等段落格式。
%
% \DescribeOption{layout}
% 设置纸张、页边距、页眉页脚和文档网格。
%
% \DescribeOption{gb7714}
% 设置参考文献后端和著录格式，见第~\ref{sec:bib-backend}~节。
%
% \DescribeOption{symbols}
% 声明符号表词条，见第~\ref{sec:abbr}~节。
%
% \DescribeOption{abbreviations}
% 声明缩略语词条，见第~\ref{sec:abbr}~节。
%
% 全文开关直接写在顶层，如 \option{openright}、\option{caption-bilingual}。
% 部件的可选参数只对该次调用生效，如 \cs{tableofcontents}\oarg{depth=3}。
% 双语内容统一写成 \texttt{\{主语言\}[副语言]}，默认主语言跟随论文的 \option{lang}：
% \begin{latex}
% \hitsetup{info={title={论文题目}[Thesis Title]}}
% \bilingualswapon  % 交换后续输入的主副语言归属
% \hitsetup{info={author={Zhang San}[张三]}}
% \bilingualswapoff
% \end{latex}
% \DescribeMacro{\bilingualswapon}
% 交换后续输入的主副语言归属，重复调用仍保持交换状态；声明遵循 \hologo{TeX} 分组。
%
% \DescribeMacro{\bilingualswapoff}
% 恢复默认输入顺序。两种声明均不修改已存数据或论文语言。
% 字段省略副值时保留副值，空值显式清空。
% \texttt{title=\{\string\keepvalue\}[New Title]} 只改副语言，\cs{keepvalue} 只作为字段保留标记，
% 不定义全局命令。章节题注是新内容，不能使用此标记。
% 非双语字段收到双语值时警告并采用主值。\texttt{report-number=\{A[B]\}} 中的方括号是普通文字。
% 值含逗号时用花括号保护，例如 \texttt{title=\{中文\}[\{Design, Analysis\}]}。
% 支持三态的键接受 \texttt{true}、\texttt{false}、\texttt{auto}；\texttt{auto} 的含义见各键说明，
% 不等于一律开启。
%
% \section{个人信息}[Document Information]
% \label{sec:titlepage}
% 以下字段写在 \cs{hitsetup}\marg{info=\{...\}} 中。
%
% \DescribeOption{title}
% 论文题目，支持双语内容及固定两行的题目栏。
%
% \DescribeOption{subtitle}
% 副题目，独立填写双语内容；修改 title 不会清空副题目。
%
% \DescribeOption{author}
% 作者姓名，可同时填写中英文姓名。
%
% \DescribeOption{supervisor}
% 主导师姓名及职称，支持双语。
%
% \DescribeOption{associate-supervisor}
% 副导师，可与主导师及行业导师同时填写，支持双语。
%
% \DescribeOption{industry-supervisor}
% 行业导师姓名及职称，支持双语。
%
% \DescribeOption{affiliation}
% 学院或所在单位，支持双语。
%
% \DescribeOption{speciality}
% 专业名称，支持双语。
%
% \DescribeOption{degree-applied}
% 申请学位全称，如“工学硕士”“Master of Engineering”。
%
% \DescribeOption{institution}
% 校名，默认哈尔滨工业大学及其英文名。
%
% \DescribeOption{student-id}
% 学号，\option{student-number} 为同义写法。
%
% \DescribeOption{secrecy}
% 论文密级，例如“公开”。
%
% \DescribeOption{classified-index}
% 国内图书分类号。
%
% \DescribeOption{udc}
% 国际图书分类号。
%
% \DescribeOption{school-code}
% 学校代码，默认 \texttt{10213}。
%
% \DescribeOption{practice-type}
% 专业学位实践成果类型，如“工程设计”。当 \texttt{degree-type=professional,form=practice} 时，
% 外封类别行显示为“（专业学位工程设计）”。不填时使用“□□……□”，值中不写括号和“专业学位”前缀。
%
% \DescribeOption{station-zh}
% 博士后流动站名称，目前仅填写中文。
%
% \DescribeOption{report-number}
% 博士后研究工作报告编号。
%
% 题目的花括号分组用于固定两行的题目栏，双反斜杠用于自由排版处换行，末尾方括号只表示副语言：
% \begin{latex}
% \hitsetup{info={title={题目第一行}{题目第二行}[English Title],
%                 subtitle={中文副题}[English Subtitle]}}
% \end{latex}
% \DescribeOption{title-en-xiaoer}
% 在 info 中同时设置外封和内封英文题目的字号：\texttt{true} 用小二，
% \texttt{false} 用二号，\texttt{auto} 按页面规则决定。
%
% \section{日期}[Dates]
% \label{sec:date}
% 日期字段也写在 info 中，采用 ISO 数字，写到哪一级就排到哪一级：
% \begin{latex}
% \hitsetup{info={date-zh={2026-06},date-en={2026-06}}}
% \end{latex}
% \DescribeOption{date-zh}
% 中文页面使用的日期，不写时使用编译当天的默认精度。
%
% \DescribeOption{date-en}
% 英文页面使用的日期，按英文日期格式输出。
%
% \DescribeOption{cover-date-zh}
% 单独覆盖中文外封日期。
%
% \DescribeOption{enroll-date-zh}
% 开题或中期报告的入学日期。
%
% \DescribeOption{research-start-date-zh}
% 博士后研究工作起始日期。
%
% \DescribeOption{research-end-date-zh}
% 博士后研究工作结束日期。
%
% \DescribeOption{report-submit-date-zh}
% 博士后报告提交日期。
%
% \DescribeOption{work-finish-date-zh}
% 博士后研究工作完成日期。页面局部参数暂不接受日期字段。
%
% \section{组织各个部件}[Organizing Document Parts]
% \label{sec:structure}
% 写哪个部件就排哪个，顺序由主文件确定。前置、正文、后置通常自动切换，也可手动控制。
%
% \DescribeMacro{\frontmatter}
% 进入前置部分，排摘要、目录等内容。
%
% \DescribeMacro{\mainmatter}
% 进入正文，开始正文页码和章节编号。
%
% \DescribeMacro{\backmatter}
% 进入后置部分，排参考文献、成果、声明等内容。
%
% \DescribeMacro{\cover}
% 排外封；开题和中期报告中排首页。
%
% \DescribeMacro{\titlepage}
% 排内封；硕士和博士还排英文题名页。
% 英文题名页放不下时，先减少字段表前的大空段，再减少题目前的大空段；
% 不因此缩小字段字号或行距。仍然放不下时给出警告，内容继续排到下一页。
%
% \DescribeEnv{abstract}
% 排一篇摘要，\texttt{[lang=en]} 选择英文摘要。
%
% \DescribeMacro{\listofsymbols}
% 根据 symbols 中的词条排符号表，也可在 \env{symbols} 环境中逐项填写。
%
% \DescribeMacro{\listofabbreviations}
% 排缩略语表。
%
% \DescribeMacro{\printnomenclature}
% 将符号和缩略语合排在同一部分。
%
% \DescribeMacro{\tableofcontents}
% 排目录，可设置语种和层级深度。
%
% \DescribeMacro{\listoffigures}
% 排图清单。
%
% \DescribeMacro{\listoftables}
% 排表清单。
%
% \DescribeMacro{\listofequations}
% 排公式清单。
%
% \DescribeEnv{conclusions}
% 排不编号的结论章。
%
% \DescribeMacro{\bibliography}
% \cs{bibliography}\marg{库名} 排参考文献。
%
% \DescribeMacro{\appendix}
% 此后的章节使用附录编号。
%
% \DescribeEnv{achievements}
% 排成果列表；博士期间的成果使用 \texttt{[period=doctoral]}。
%
% \DescribeMacro{\defense}
% 排答辩页，支持填写名单与决议正文或插入扫描件。
%
% \DescribeMacro{\declarations}
% 排原创性声明与使用授权页。
%
% \DescribeEnv{acknowledgements}
% 排致谢正文。
%
% \DescribeEnv{resume}
% 排个人简历。
%
% \DescribeEnv{permanentaddress}
% 排博士后永久通讯地址。
%
% \DescribeMacro{\printindex}
% 排索引。
%
% \DescribeOption{shown}
% 部件显示开关，一般接受 \texttt{auto}、\texttt{true}、\texttt{false}。
% \texttt{auto} 按阶段与学位决定；报告里终稿专有部件默认不排。部分报告部件不能强制排出，
% 写 \texttt{shown=true} 时会提示。部件通常只排一次，重复调用会提示。
%
% \section{外封与内封的局部设置}[Local Cover and Title Page Settings]
% \cs{cover} 和 \cs{titlepage} 接收本页的元信息覆盖，排完恢复全文值：
% \begin{latex}
% \cover[info={title={外封题目}[Cover Title],author={张三}[Zhang San]}]
% \titlepage[info={title={内封题目}[English Title],
%                 subtitle={副标题}[Subtitle],school-code={10213}}]
% \end{latex}
% 可覆盖题目、副题目、作者、三类导师、院系、专业、申请学位、校名，以及学号、密级、分类号等。
% 全文元信息和 PDF 属性不随本页覆盖改变。\env{titlepage} 环境仍可排自定义整页。
%
% \DescribeOption{show-subtitle}
% 副题目显示开关：\texttt{auto} 按当前语言副题目是否为空决定显示，\texttt{true/false} 强制显示或隐藏。
% 题目和副题目分别更新，互不清空。
% 使用 \texttt{\string\hitsetup\{show-subtitle=false\}} 隐藏副题目。
%
% \DescribeOption{degree-type-line}
% 外封的学位类别行显示开关，写在 \cs{cover} 的可选参数中。
% 外封还可局部覆盖 \option{practice-type}。
%
% \section{摘要与目录}[Abstracts and Contents]
% \DescribeOption{keywords}
% 摘要关键词用半角逗号分隔，在环境参数中按该篇摘要语言填写；省略时读取 \option{info/keywords} 的对应语言。
%
% \DescribeOption{keywords-hanging}
% 设为 \texttt{true} 时，关键词续行对齐关键词正文。
% \begin{latex}
% \begin{abstract}[lang=zh,keywords={有限元,优化},keywords-hanging=true]
% 摘要正文。
% \end{abstract}
% \tableofcontents[depth=3,lang={zh,en}]
% \end{latex}
% \DescribeOption{depth}
% 目录深度，按级数计数，默认 3，写在 \cs{tableofcontents} 的可选参数中。
%
% \DescribeOption{lang}
% 目录语种：\texttt{auto} 按论文类型决定，\texttt{zh/en} 选择一种语言，\texttt{\{zh,en\}} 排两份目录。
% 报告只排一份目录，语种跟随文档。英文目录条目来自标题末尾参数，如
% \cs{chapter}\marg{绪论}\oarg{Introduction}；缺英文标题时请按日志提示补齐。
%
% \DescribeOption{latin-sans}
% 控制目录条目的西文是否使用无衬线字体。
%
% \DescribeOption{latin-natural}
% 控制纯西文目录条目是否按西文字体的自然行高排版。
%
% \DescribeOption{lists}
% 控制图表清单是否列入目录。
%
% \DescribeOption{front-blank}
% 控制目录中第一个正文章节条目之前是否空一行。
%
% \section{一份配置用于多个变体}[Conditional Configuration]
% \DescribeMacro{\hitif}
% \cs{hitif}\oarg{条件}\marg{成立时}\marg{不成立时} 按条件选择内容；\cs{hitsetup} 也接受条件参数。
% \begin{latex}
% \hitsetup[degree-level=master,campus=harbin]{info={student-id={1234567890}}}
% \hitif[stage=final]{终稿内容}{报告内容}
% \end{latex}
% 条件中逗号表示“且”，分号表示“或”，同一项的多个值用竖线分隔，叹号表示取反。
%
% \chapter{正文写作}[Writing the Main Text]
% \section{标题、分页与编号}[Headings, Page Breaks, and Numbering]
% \DescribeMacro{\chapter}
% 章标题，写作 \cs{chapter}\marg{主语言标题}\oarg{副语言标题}。
% 报告可从章或节开始，模板按首个编号标题确定映射；确定后不要混用两套层级。
%
% \DescribeMacro{\section}
% 节标题，双语参数写法同 \cs{chapter}。
%
% \DescribeMacro{\subsection}
% 条标题，双语参数写法同 \cs{chapter}。
%
% \DescribeMacro{\subsubsection}
% 款标题，双语参数写法同 \cs{chapter}。
% \begin{latex}
% \chapter{研究方法}[Methods]
% \section{模型}[Model]
% \hitsetup{
%   top-level-heading-pagebreak-before=false,
%   caption-numbering-by-top-level-heading=false,
%   equation-numbering-by-top-level-heading=true,
% }
% \end{latex}
% 顶层在终稿中是章，在报告中按标题映射确定。
%
% \DescribeOption{top-level-heading-pagebreak-before}
% 顶层标题前是否换页：\texttt{true} 换页，\texttt{false} 连续排版，默认 \texttt{auto} 沿用变体规则。
% 开关为 true 或模板默认要求换页时，先排完待处理图表，再换页并连续分节；
% 终稿沿用当前右开页设置。新节首段按 Word 规则保留页首段前，不再受兼容模式 14/15 的差别影响。
% 段前与上一段段后的合并或相加规则仍然有效；连续标题不一定保留段前参数的全额。
%
% \DescribeMacro{\sectionbreak}
% 插入 Word 意义的分节符，与章节标题 \cs{section} 不同，不改变章节编号。
% 类型使用裸选项 \texttt{continuous/next/odd/even}，一次只能选一种；
% 不写类型或写 \texttt{auto} 时取 \texttt{next}。
% continuous 结束当前段落但不主动换页；next 清完浮动图表后从下一页开始，
% odd、even 分别从奇数页、偶数页开始，必要时补不印页眉页脚的空白页。
% \begin{latex}
% \sectionbreak[continuous]
% \sectionbreak[next]
% \sectionbreak[odd]
% \sectionbreak[even]
% \end{latex}
% 新节默认继承设置并连续编页码；按节重编号的脚注会重新编号。
% 同一组选项还可传入 docxlike 的 columns、page-setup 和 page-number-format，
% 这三组取 auto 表示不覆盖；子键沿用 docxlike 的名称。
% \begin{latex}
% \sectionbreak[continuous,
%   columns={number-of-columns=2,spacing=1cm}]
% \sectionbreak[continuous,columns={number-of-columns=1}]
% \sectionbreak[next,
%   page-number-format={number-format=lower-roman,start-at=1}]
% \end{latex}
% 连续分节不能修改纸张、页边距、网格或页眉页脚距边界；需要改这些尺寸或重启页码时用分页型分节。
% 页中连续分节的本页页眉页脚仍取前一节，空白页首分节则立即使用新节。
% 模板的页眉页脚内容仍由页面样式控制；底层节属性不会自动改写模板的页眉文字。
% 原生 \cs{newpage}、\cs{clearpage}、\cs{cleardoublepage} 保持 LaTeX 的含义，不自动分节。
% 分章应清完上一章的浮动体；手动控制时可写：
% \begin{latex}
% \clearpage
% \sectionbreak[continuous]
% \end{latex}
%
% \DescribeOption{caption-numbering-by-top-level-heading}
% 图表、算法和代码清单是否按顶层标题编号，支持 \texttt{auto}（默认）、\texttt{true}、\texttt{false}。
% 开启时形如“2-1”，关闭后全文连续编号。
%
% \DescribeOption{equation-numbering-by-top-level-heading}
% 公式是否按顶层标题编号，支持 \texttt{auto}（默认）、\texttt{true}、\texttt{false}，与题注编号独立。
%
% \DescribeOption{openright}
% 是否补空白页以保证右开页；\texttt{false} 取消补页，不等于取消标题换页。
%
% \DescribeOption{equation-numbering-fullwidth}
% 控制公式号括号的全半角；英文文档始终使用半角。
%
% \section{图、表与分图}[Figures, Tables, and Subfigures]
% \DescribeMacro{\caption}
% 单语题注使用标准 \cs{caption}；双语题注使用 \texttt{\{主语言\}[副语言]}。
% 中文文档默认先中文、后英文：
% \begin{latex}
% \begin{figure}[htbp]
%   \centering
%   \includegraphics[width=.7\linewidth]{figures/example}
%   \caption{中文图题}[English caption]
%   \label{fig:example}
% \end{figure}
% \end{latex}
% \cs{label} 放在题注后，正文用 \cs{ref} 引用。
%
% \DescribeOption{caption-bilingual}
% 控制是否排副语言题注，默认按学位决定。
%
% 表题写在表格上方，三线表可用 \pkg{booktabs} 的 \cs{toprule}、\cs{midrule}、\cs{bottomrule}。
% 跨页表格使用 \pkg{longtable} 或 \pkg{tabularray} 的 \env{longtblr}，续图表使用 \cs{ContinuedFloat}。
% \DescribeMacro{\subcaptionbox}
% 分图推荐使用 \pkg{subcaption} 提供的 \cs{subcaptionbox}：
% \begin{latex}
% \begin{figure}[htbp]
%   \centering
%   \subcaptionbox{{子图甲\label{fig:a}}[Part A]}[.45\linewidth]{\rule{3cm}{2cm}}
%   \hfill
%   \subcaptionbox{{子图乙\label{fig:b}}[Part B]}[.45\linewidth]{\rule{3cm}{2cm}}
%   \caption{总图题}[Overall caption]
% \end{figure}
% \end{latex}
% 细调分图使用 \cs{captionsetup}\oarg{sub}\marg{...}；模板尚未提供 iota 的四个
% \texttt{subcaption-*} 顶层键。图表前后的默认空段由模板处理，不要重复添加空行。
% 也可使用 \env{subfigure} 环境，在其中写 \cs{caption}\marg{主语言}\oarg{副语言}。
% \env{longtable} 内的题注使用同一命令；\env{longtblr} 的外部选项写作
% \texttt{caption=\{\{中文表题\}[English Caption]\}}。
% 英文文档默认先英文、后中文；\cs{bilingualswapon} 可局部交换输入顺序。
% 前置可选参数仍是文档语言的短题注，例如 \cs{caption}\oarg{短图题}\marg{完整中文图题}\oarg{English Caption}，
% 不随交换声明改变语言。
% 示例展示三种插图方式：四张轴承矢量 PDF 可用 \cs{includegraphics} 插入；
% 第 1 章第一幅静压轴承图直接在正文写 \env{tikzpicture}；第 1 章及深圳本科中期报告使用
% \cs{input}\marg{figures/volute-section} 引入外部 Ti\emph{k}Z 源码。
% 两种 Ti\emph{k}Z 写法均在论文编译时绘图，导言区须加载 \pkg{tikz} 及其
% \texttt{arrows.meta} 库，示例主文件已配置。
% 资源源码见 \file{src/assets/hit-example-figures.dtx}，修改后运行 \texttt{make distribute}
% 生成四张 PDF，并把它们和涡室截面的 \file{.tex} 文件复制到示例的 \file{figures/}。
% 无编号题注使用 \cs{caption*}\marg{中文说明}\oarg{English Note}：开启双语题注时排出两种文字，
% 不加图表编号，也不生成清单条目。在 \env{longtable} 中，题注后还需写 \verb|\\| 结束该行；
% 这里的星号不表示表下注释或自动续表标题。
%
% \section{公式与定理}[Equations and Theorems]
% \label{sec:math}
% 使用标准 \env{equation}、\env{align} 等数学环境和 \cs{eqref} 引用。
% \DescribeEnv{eqdenote}
% 对齐公式后的符号说明：
% \begin{latex}
% \begin{equation}\label{eq:force}F=ma\end{equation}
% \begin{eqdenote}
%   \item $F$ & 力，N；
%   \item $m$ & 质量，kg；
%   \item $a$ & 加速度，m/s$^2$。
% \end{eqdenote}
% \end{latex}
% \env{eqdenote}\oarg{其中} 可改引导词，空参数不印引导词。
% 定理环境包括 \env{theorem}、\env{lemma}、\env{definition}、\env{proposition}、
% \env{corollary}、\env{example}、\env{remark} 和 \env{proof} 等。
%
% \DescribeOption{theorem-numbering-by-top-level-heading}
% 控制定理编号是否随顶层标题重置。
%
% \DescribeOption{theorem-body-indent}
% 设置定理头与正文之间的间隔；更多样式写在 \option{thm} 族中。
%
% \section{列表与代码}[Lists and Code]
% \env{enumerate} 是编号列表，\env{itemize} 是项目列表，\env{description} 是说明列表。
% \begin{latex}
% \hitsetup{enumerate-hanging=true,itemize-hanging=false}
% \begin{enumerate}
%   \item 第一项。
%   \item 第二项。
% \end{enumerate}
% \end{latex}
% \DescribeOption{enumerate-hanging}
% 编号列表的续行对齐：开启时对齐首行正文，关闭时回到该级列表左缘，\texttt{auto} 保留变体默认。
%
% \DescribeOption{itemize-hanging}
% 项目列表的续行对齐，取值含义与 \option{enumerate-hanging} 相同。
% 环境可选参数中显式指定的 \pkg{enumitem} 尺寸优先。
%
% \DescribeEnv{lstlisting}
% 使用 \pkg{listings} 排代码清单，行内代码可使用 \cs{verb}。
%
% \DescribeEnv{algorithm}
% 使用 \pkg{algorithm2e} 排伪代码，默认已加载相应宏包。
%
% \DescribeOption{raw-style}
% 设置代码清单的外框和行号样式。
%
% \DescribeOption{algorithm-style}
% 设置伪代码的外框、关键字大小写、缩进引导线和行号样式。
% 算法使用同一 \cs{caption} 命令；代码清单通过键值参数给题注：
% \begin{latex}
% \begin{lstlisting}[language=Python,caption={{代码示例}[Code Example]}]
% print(1)
% \end{lstlisting}
% \end{latex}
% \DescribeOption{verb-space}
% \texttt{cell} 隐藏行内代码的空格记号，\texttt{visible} 显示它。
%
% \DescribeOption{code-overhang}
% 代码断行增强开关，取 \texttt{true/false/auto}；默认 \texttt{auto}，按 \texttt{true} 处理。
% 适用于 \hologo{LuaLaTeX} 下的 \cs{verb}、\cs{texttt} 和局部 \cs{ttfamily} 文本。
% 默认使用贪心断行，保留原生及用户显式断点，并补充逗号、单个赋值等号、连字符和路径分隔符附近的软断点。
% \texttt{false} 关闭自动补点；实际是否允许伸出版心由下项安全距离决定，默认不伸出。
% 只有实际行尾为代码时才允许扩展，代码后仍有普通文字的行
% 保持原来的版心与对齐。伸出版心的行尾代码保留自然间距，长串优先在已有断点折行，
% 单个长串仍放不下时按字形折行。列表、分栏和盒子内不借用容器外的空间。
% \cs{texttt} 仍解析 \hologo{TeX} 命令，特殊字符需要转义；\cs{verb} 按原样读取字符。
%
% \DescribeOption{code-overhang-page-inset}
% 从页面右边缘向内保留的安全距离，默认及 \texttt{auto} 均为 \texttt{100 percent}，
% 即保留全部正文右页边空白，代码在版心内折行。改为 \texttt{50 percent} 时最多借用一半右页边空白。
% 百分比可取 0 到 100；也可指定 \texttt{10mm} 等长度。
% 它不是超出版心的长度，也不修改正文页边距。双面文档取奇偶页都能容纳的宽度。
% 手册与论文、报告采用同一组默认值。\pkg{listings} 等自行断行的代码环境仍使用各自的折行设置，
% 严格字符网格模式仍按网格断行。
% \begin{latex}
% % 恢复模板默认：贪心、增强断点、不伸出版心
% \hitsetup{code-overhang=auto,code-overhang-page-inset=auto}
% % 显式允许借用一半右页边空白
% \hitsetup{code-overhang=true,code-overhang-page-inset=50 percent}
% \end{latex}
%
% \DescribeMacro{\hitenter}
% 必须增加空段时使用 \cs{hitenter}\oarg{行数}；需要在页首保留空段时用 \cs{hitenter*}。
%
% \chapter{符号、文献与后置材料}[Symbols, References, and Back Matter]
% \section{符号与缩略语}[Symbols and Abbreviations]
% \label{sec:abbr}
% \begin{latex}
% \hitsetup{
%   symbols={ {$p$}={压力，Pa}[{Pressure, Pa}], {$v$}={速度}[Velocity] },
%   abbreviations={ FEM={long={有限元法}[Finite Element Method],
%     long-plural={\keepvalue}[Finite Element Methods]} },
% }
% \end{latex}
% 符号按输入顺序排，也可用组名嵌套一层分组：
% \begin{latex}
% \hitsetup{symbols={
%   拉丁字母={
%     {$p$}={压力}[Pressure],
%     {$v$}={速度}[Velocity],
%   },
%   希腊字母={ {$\rho$}={密度}[Density] },
% }}
% \listofsymbols
% \end{latex}
% 上例在同一张符号表内插入“拉丁字母”“希腊字母”两个分组标题，组内说明按当前语言选取。
% 分组标题目前为单份文字，不随语言自动切换；分组只支持一层。
%
% \DescribeEnv{symbols}
% 允许直接写 \cs{item}\oarg{符号}及说明。
%
% \DescribeOption{labelwidth}
% 指定符号表和缩略语表的标签列宽，省略时按内容计算。
%
% \DescribeMacro{\abbr}
% \cs{abbr}\marg{FEM} 首次展开全称，以后排简称；星号形式强制全称且不改变首次使用状态。
%
% \DescribeMacro{\abbrs}
% 输出缩略语复数形式，首次展开与星号行为同 \cs{abbr}。
%
% \DescribeOption{long}
% 缩略语全称，使用 \texttt{\{主语言\}[副语言]} 填写。
%
% \DescribeOption{short}
% 缩略语的简称，省略时使用词条标识符。
%
% \DescribeOption{plural}
% 简称的复数形式，省略时在简称后加 \texttt{s}。
%
% \DescribeOption{long-plural}
% 双语全称的复数形式；中文默认沿用全称，英文默认加 \texttt{s}。
%
% \DescribeOption{indexed}
% 是否为当前缩略语登记索引，可覆盖全文设置。
%
% \DescribeOption{pinyin}
% 中文索引的排序读音，如 \texttt{zhong4}。
%
% 使用 \cs{listofsymbols} 和 \cs{listofabbreviations} 分别排表，或只写 \cs{printnomenclature} 合排。
% 缩略语表依赖使用记录，新增词条后交给 latexmk 完成多轮编译。
%
% \DescribeOption{abbreviation-fullwidth}
% 控制首次展开时括号和逗号的全半角，\texttt{auto} 跟随语种。
%
% \DescribeOption{abbreviation-indexed}
% 是否为缩略语同时登记索引，可由词条的 indexed 覆盖。
%
% \DescribeOption{abbreviation-links}
% \texttt{true/auto} 允许链接到已经排出的对应表项，\texttt{false} 关闭链接。
%
% \DescribeOption{abbreviation-reset-after-abstract}
% 控制每篇摘要之后是否重置缩略语的首次使用状态。
%
% \DescribeMacro{\hitindex}
% \cs{hitindex}\marg{词} 输出文字并登记索引，最后用 \cs{printindex} 排索引。
% 中文多音字可写 \texttt{\string\hitindex[pinyin=zhong4]\{重\}} 指定排序读音。
%
% \section{参考文献}[References]
% \label{sec:bib-backend}
% 默认使用 \hologo{BibTeX}。
%
% \DescribeMacro{\cite}
% \cs{cite}\marg{key} 在正文中生成上标引用。
%
% \DescribeMacro{\inlinecite}
% \cs{inlinecite}\marg{key} 在正文中生成行内引用。
%
% 文献表写 \cs{bibliography}\marg{references}，库文件为 \file{references.bib}。
% 使用 \hologo{biber} 时，在导言区设置后端并登记文献库：
% \begin{latex}
% \hitsetup{gb7714={backend=biber}}
% \addbibresource{references.bib}
% \end{latex}
% 正文仍可写 \cs{bibliography}\marg{references}，或使用 \cs{printbibliography}。
% latexmk 自动选择 \hologo{BibTeX} 或 \hologo{biber}。\option{backend=biber} 不可放到正文。
% 以下键写在 gb7714 中。
%
% \DescribeOption{backend}
% 参考文献后端：\texttt{bibtex} 使用 \hologo{BibTeX}，\texttt{biber} 使用 \pkg{biblatex} 与 \hologo{biber}。
% 默认 \texttt{auto}，未自行加载 \pkg{biblatex} 时使用 \hologo{BibTeX}。
%
% \DescribeOption{compress-min}
% 连续引用从 2 个或 3 个开始压缩为区间。
%
% \DescribeOption{uppercase-family}
% 控制著者姓氏是否全部大写。
%
% \DescribeOption{max-bib-names}
% 参考文献表中显示的作者数上限。
%
% \DescribeOption{max-cite-names}
% 引用中显示的作者数上限。
%
% \DescribeOption{punct}
% 著录标点：\texttt{full} 使用全角，\texttt{half} 使用半角，\texttt{by-lang} 按条目语言选择。
%
% \DescribeOption{sentence-case}
% 控制题名是否转为句首大写，可分别设置 title、booktitle、journal。
%
% \DescribeOption{numbering-align}
% 控制文献标号在标号列中的对齐方式。
%
% \DescribeOption{split-entry}
% 控制单条参考文献是否允许跨页。
%
% \DescribeOption{show-url}
% 控制是否显示文献网址。
%
% \DescribeOption{show-doi}
% 控制是否显示 DOI。
%
% \DescribeOption{show-note}
% 控制是否显示条目末尾的附注。
%
% \DescribeOption{bibtex/style}
% \hologo{BibTeX} 样式名，默认 \texttt{hithesis}。
%
% \DescribeOption{biber}
% 传给 \pkg{biblatex} 的加载选项，应与后端一起设置；两种后端的著录细节可能不同。
%
% \section{成果、答辩与声明}[Achievements, Defense, and Declarations]
% \DescribeEnv{pubgroup}
% 成果推荐使用单独的文献库，并按组选择条目：
% \begin{latex}
% \begin{achievements}[publications]
%   \begin{pubgroup}{发表的学术论文}
%     \item{paper-one}
%     \item{paper-two}
%   \end{pubgroup}
% \end{achievements}
% \end{latex}
% \hologo{BibTeX} 需要额外处理成果库，示例 \file{latexmkrc} 已包含此步骤。
% \hologo{biber} 用户在导言区用 \cs{addbibresource} 登记成果库，环境上不再写库名。
% 答辩名单是 \cs{defense} 的局部参数：
% \begin{latex}
% \defense[
%   unframed=true,
%   reviewers={{name={张三},title={教授},affiliation={某大学},discipline={某学科}}},
%   chair={name={李四},title={教授},affiliation={某大学}},
%   members={{name={王五},title={教授}}},
%   secretary={name={赵六}},
%   resolution={答辩委员会决议正文。},
% ]
% \declarations
% \end{latex}
% 以下键写在 \cs{defense} 的可选参数中。
%
% \DescribeOption{unframed}
% 取 \texttt{true}、\texttt{false} 或 \texttt{auto}；裸写 \texttt{unframed} 等同于 \texttt{true}。
% \texttt{auto}（默认）沿用当前有框兼容版，\texttt{false} 显式选择有框版。
% \texttt{true} 生成无框名单，只排实际填写的人员，名单可跨页并重复表头，决议按普通正文接排。
% 不填写 \option{reviewers} 时不生成评阅人记录，适用于匿名评审；无框名单不设签名栏。
% 2025 年写作指南要求博士增加这一部分，并规定采用无框线表格，但学校现有书写范例仍为有框样表。
% 此可选模式依据指南文字实现；列宽、留白及分页处理是本模板的实现选择，不是对学校新版样表的复刻。
% \option{file} 插入的扫描件不会因本开关重排。无框标签可用 \texttt{overrides/defense} 中的
% \texttt{unframed-reviewer}、\texttt{unframed-column-title}、\texttt{unframed-text} 等词条覆盖。
%
% \DescribeOption{reviewers}
% 评阅人记录列表，每条可填写 name、title、affiliation、discipline。
%
% \DescribeOption{chair}
% 答辩委员会主席记录，字段写法同评阅人。
%
% \DescribeOption{members}
% 答辩委员会成员记录列表。
%
% \DescribeOption{secretary}
% 答辩秘书记录。
%
% \DescribeOption{resolution}
% 答辩委员会决议正文。
%
% \DescribeOption{file}
% 在 \cs{defense} 或 \cs{declarations} 中指定签字后的扫描件，例如 \texttt{file=defense.pdf}。
%
% \DescribeOption{page-number}
% 设为 \texttt{true} 时，在插入的扫描件页面上显示模板页码。
%
% 致谢和简历直接写入同名环境，附录先写 \cs{appendix}，再用常规章节命令。
%
% \chapter{调整与排错}[Customization and Troubleshooting]
% \section{改样式与固定文字}[Customizing Styles and Terms]
% \label{sec:const}
% 不要直接修改生成的 \file{hithesis.cls} 或 \file{hithesis.cfg}，在导言区覆盖即可。
% \begin{latex}
% \hitsetup{
%   styles={Heading 1={font={font-style=regular}}},
%   overrides={titlepage={supervisor={指导教师}[Supervisor]},
%              references={title={参考文献}[References]}},
% }
% \end{latex}
% \option{styles} 按 Word 样式组织，如 \texttt{Normal}、\texttt{Heading 1}、\texttt{Header}、\texttt{Footer}，
% 使用 \option{font}、\option{paragraph} 等子组。页面设置放在 \option{layout} 中。
% \option{overrides} 按页面分组，词条使用相同的双语输入格式；常用 \texttt{cover}、\texttt{titlepage}、\texttt{abstract}、
% \texttt{references}、\texttt{achievements}、\texttt{defense}、\texttt{declarations}；\texttt{base} 为跨页共用文字。
% 字段、样式和词条是三类设置：改作者姓名用 \option{info}，改字号用 \option{styles}，改“作者”标签用 \option{overrides}。
% 图表名称放在 \texttt{caption} 中，部件标题统一使用 \texttt{title}。例如：
% \begin{latex}
% \hitsetup{overrides={
%   caption={figure={照片}[Photograph], equation={式}[Equation]},
%   list-of-figures={title={照片清单}[List of Photographs]},
%   list-of-equations={title={公式清单}[List of Equations],
%                      equation={公式}[Equation]},
% }}
% \end{latex}
% \texttt{caption/equation} 是正文引用的名称，默认“式”；
% \texttt{list-of-equations/equation} 是公式清单的条目名称，默认“公式”。
% 图表清单的类型名分别为 \texttt{list-of-figures/image} 和 \texttt{list-of-tables/table}，
% 默认跟随题注名称，也可单独覆盖。
% 图、表和公式的 \cs{autoref} 名称随这些词条更新。公式编号的括号仍由
% \option{equation-numbering-fullwidth} 设置，不写进名称里。
% 标准名称只保存当前文档所用的语言；单个值直接用于当前语言，双语值按正文语言选取。
% 原来的 \texttt{figurename}、\texttt{tablename} 等写法仍可使用，与分组写法更新同一处取值。
% 算法关键字归入 \texttt{overrides/algorithm}，学位类型文字归入
% \texttt{overrides/degree-type}；编号后的内容可写成
% \texttt{overrides=\{number-separator=\{section=\{\textbackslash quad\}\}\}}。
% 封面上的完整文种统一使用 \texttt{overrides/cover/document-type}，例如“硕士学位论文”或
% “硕士学位论文开题报告”；\texttt{overrides/cover/title} 是“题目：”这类标签，
% 用户填写的研究题目仍使用 \texttt{info/title}。三者分别设置，排版时不再给文种追加学位或阶段。
% 原报告标题和题目标签专用键已移除，不保留转发入口。
% 校区等档位写在其他字段名末尾，例如 \texttt{cover/report-student-id-shenzhen}
% 和 \texttt{cover/report-imprint-shenzhen}。
% \texttt{defense/title} 是答辩页标题，职称列表头使用 \texttt{defense/column-title}。
% 需要完整的底层键时查 docxlike 文档及仓库配置文件，不要套用 Typst 的键值语法。
% 条件后缀可按任意顺序输入，内部按阶段、校区、学位、学位类型、学科门类、成果形式排序。
% 同一轴不能重复指定；同一组 \texttt{overrides} 中规范化后重名也会报错。
%
% \subsection{可选的智能引用}
% wiseref 是独立宏包，需要用户自行安装并加载，模板不会自动加载它。
% \begin{latex}
% \usepackage{wiseref}
% 参见\wref{fig:structure,tab:result}。
% \wref[照片]{fig:structure}
% \wref[]{eq:balance}[]
% \wref*{fig:structure}
% \end{latex}
% \cs{wref} 的标签名不限写法；两个可选参数分别覆盖前、后缀，省略时沿用设置，
% 空参数清除对应部分。星号形式不带链接。编号格式独立于前后缀，因此公式引用
% 清除名称后仍保留括号，括号跟随 \option{equation-numbering-fullwidth}。
% 加载宏包后，模板根据现有词条设置名称；用户随后可用 \cs{wiserefsetup} 覆盖。
% 语言选择使用 \texttt{lang}，默认跟随文档设置；显式格式不会因语言选择变化而重置。
% \begin{latex}
% \wiserefsetup{figure={prefix={照片}},sort=false,compress=false}
% \end{latex}
% 默认按输入顺序排列，相邻同类引用合组。排序和连续编号压缩需显式开启，
% 自定义公式号与字母子编号不按普通数字压缩。详细配置见 wiseref 的说明文档。
%
% \section{旧文档迁移}[Migrating Existing Documents]
% \label{sec:upgrade}
% 先备份论文，再换模板，并根据日志逐项迁移。
%
% \DescribeMacro{\documentclass}
% 旧类名 hithesisbook、hithesisart、hithesisartplus 改为 \hithesis\ 的统一文档类 \cls{hithesis}，
% 配合阶段、学位与校区选项。type 改为 degree-level，opening/midterm 改为 proposal/interim；
% 本科毕业设计写作 \texttt{degree-level=bachelor,form=practice}，不再使用 \texttt{bachelor*}。
%
% \DescribeMacro{\makecover}
% 改为分别调用 \cs{cover}、\cs{titlepage} 和摘要环境。
%
% \DescribeMacro{\maketitle}
% 改为 \cs{titlepage}。
%
% \DescribeOption{term}
% 固定词条配置改用 overrides，双语值写成 \texttt{\{主语言\}[副语言]}。
%
% \DescribeOption{const}
% 与旧 term 配置相同，改用 overrides。
%
% \DescribeOption{bib}
% 参考文献配置改用 gb7714。
%
% \DescribeOption{toc}
% 旧目录配置族改为 \cs{tableofcontents} 的局部参数。
%
% \DescribeOption{defense}
% 旧答辩配置族改为 \cs{defense} 的局部参数。
%
% \DescribeOption{co-supervisor}
% 按实际角色改用 associate-supervisor（副导师）或 industry-supervisor（行业导师）。
%
% 已发布版本的部分旧命令仍有兼容提示，开发版中未发布的配置族可能直接报错。
% 新文档请从当前示例开始，不混用不同版本的类文件和配置文件。
%
% \section{常见问题}[Frequently Asked Questions]
% \label{sec:howtoask}
% \begin{description}
% \item[找不到字体] 检查字体是否安装、\option{fontset} 是否对应当前系统；可先用 \texttt{fandol} 确认能编译。
% \item[参考文献或索引是空的] 使用项目自带 \file{latexmkrc} 编译，检查库路径、引用键以及 \hologo{BibTeX}/\hologo{biber} 日志。
% \item[出现问号或英文目录错误提示] 补齐标签和英文标题，让 latexmk 完成后续编译轮次。
% \item[报告中没有某个整页部件] 检查该部件是否适用于当前 \option{stage} 和 \option{degree-level}。
% \item[封面内容溢出] 缩短过长字段或调整题目分行，英文题目可设置 \option{title-en-xiaoer=true}。
% \item[升级后输出没有变化] 检查实际加载的类文件路径，清理旧辅助文件后重编，避免使用另一目录的旧副本。
% \end{description}
% 提交问题时附最小可复现的 \file{.tex}、编译日志、\hologo{TeX}~Live 版本、引擎和类选项。
% 源码修改、构建与测试说明见 \file{CONTRIBUTING.md}；版本变化见本手册的更新日志。
%
% \Finale
